\IfFileExists{stacks-project.cls}{%
\documentclass{stacks-project}
}{%
\documentclass{amsart}
}

% For dealing with references we use the comment environment
\usepackage{verbatim}
\newenvironment{reference}{\comment}{\endcomment}
%\newenvironment{reference}{}{}
\newenvironment{slogan}{\comment}{\endcomment}
\newenvironment{history}{\comment}{\endcomment}

% For commutative diagrams we use Xy-pic
\usepackage[all]{xy}

% We use 2cell for 2-commutative diagrams.
\xyoption{2cell}
\UseAllTwocells

% We use multicol for the list of chapters between chapters
\usepackage{multicol}

% This is generally recommended for better output
\usepackage{lmodern}
\usepackage[T1]{fontenc}

% For cross-file-references
\usepackage{xr-hyper}

% Package for hypertext links:
\usepackage{hyperref}

% For any local file, say "hello.tex" you want to link to please
% use \externaldocument[hello-]{hello}
\externaldocument[introduction-]{introduction}
\externaldocument[conventions-]{conventions}
\externaldocument[sets-]{sets}
\externaldocument[categories-]{categories}
\externaldocument[topology-]{topology}
\externaldocument[sheaves-]{sheaves}
\externaldocument[sites-]{sites}
\externaldocument[stacks-]{stacks}
\externaldocument[fields-]{fields}
\externaldocument[algebra-]{algebra}
\externaldocument[brauer-]{brauer}
\externaldocument[homology-]{homology}
\externaldocument[derived-]{derived}
\externaldocument[simplicial-]{simplicial}
\externaldocument[more-algebra-]{more-algebra}
\externaldocument[smoothing-]{smoothing}
\externaldocument[modules-]{modules}
\externaldocument[sites-modules-]{sites-modules}
\externaldocument[injectives-]{injectives}
\externaldocument[cohomology-]{cohomology}
\externaldocument[sites-cohomology-]{sites-cohomology}
\externaldocument[dga-]{dga}
\externaldocument[dpa-]{dpa}
\externaldocument[sdga-]{sdga}
\externaldocument[hypercovering-]{hypercovering}
\externaldocument[schemes-]{schemes}
\externaldocument[constructions-]{constructions}
\externaldocument[properties-]{properties}
\externaldocument[morphisms-]{morphisms}
\externaldocument[coherent-]{coherent}
\externaldocument[divisors-]{divisors}
\externaldocument[limits-]{limits}
\externaldocument[varieties-]{varieties}
\externaldocument[topologies-]{topologies}
\externaldocument[descent-]{descent}
\externaldocument[perfect-]{perfect}
\externaldocument[more-morphisms-]{more-morphisms}
\externaldocument[flat-]{flat}
\externaldocument[groupoids-]{groupoids}
\externaldocument[more-groupoids-]{more-groupoids}
\externaldocument[etale-]{etale}
\externaldocument[chow-]{chow}
\externaldocument[intersection-]{intersection}
\externaldocument[pic-]{pic}
\externaldocument[weil-]{weil}
\externaldocument[adequate-]{adequate}
\externaldocument[dualizing-]{dualizing}
\externaldocument[duality-]{duality}
\externaldocument[discriminant-]{discriminant}
\externaldocument[derham-]{derham}
\externaldocument[local-cohomology-]{local-cohomology}
\externaldocument[algebraization-]{algebraization}
\externaldocument[curves-]{curves}
\externaldocument[resolve-]{resolve}
\externaldocument[models-]{models}
\externaldocument[functors-]{functors}
\externaldocument[equiv-]{equiv}
\externaldocument[pione-]{pione}
\externaldocument[etale-cohomology-]{etale-cohomology}
\externaldocument[proetale-]{proetale}
\externaldocument[relative-cycles-]{relative-cycles}
\externaldocument[more-etale-]{more-etale}
\externaldocument[trace-]{trace}
\externaldocument[crystalline-]{crystalline}
\externaldocument[spaces-]{spaces}
\externaldocument[spaces-properties-]{spaces-properties}
\externaldocument[spaces-morphisms-]{spaces-morphisms}
\externaldocument[decent-spaces-]{decent-spaces}
\externaldocument[spaces-cohomology-]{spaces-cohomology}
\externaldocument[spaces-limits-]{spaces-limits}
\externaldocument[spaces-divisors-]{spaces-divisors}
\externaldocument[spaces-over-fields-]{spaces-over-fields}
\externaldocument[spaces-topologies-]{spaces-topologies}
\externaldocument[spaces-descent-]{spaces-descent}
\externaldocument[spaces-perfect-]{spaces-perfect}
\externaldocument[spaces-more-morphisms-]{spaces-more-morphisms}
\externaldocument[spaces-flat-]{spaces-flat}
\externaldocument[spaces-groupoids-]{spaces-groupoids}
\externaldocument[spaces-more-groupoids-]{spaces-more-groupoids}
\externaldocument[bootstrap-]{bootstrap}
\externaldocument[spaces-pushouts-]{spaces-pushouts}
\externaldocument[spaces-chow-]{spaces-chow}
\externaldocument[groupoids-quotients-]{groupoids-quotients}
\externaldocument[spaces-more-cohomology-]{spaces-more-cohomology}
\externaldocument[spaces-simplicial-]{spaces-simplicial}
\externaldocument[spaces-duality-]{spaces-duality}
\externaldocument[formal-spaces-]{formal-spaces}
\externaldocument[restricted-]{restricted}
\externaldocument[spaces-resolve-]{spaces-resolve}
\externaldocument[formal-defos-]{formal-defos}
\externaldocument[defos-]{defos}
\externaldocument[cotangent-]{cotangent}
\externaldocument[examples-defos-]{examples-defos}
\externaldocument[algebraic-]{algebraic}
\externaldocument[examples-stacks-]{examples-stacks}
\externaldocument[stacks-sheaves-]{stacks-sheaves}
\externaldocument[criteria-]{criteria}
\externaldocument[artin-]{artin}
\externaldocument[quot-]{quot}
\externaldocument[stacks-properties-]{stacks-properties}
\externaldocument[stacks-morphisms-]{stacks-morphisms}
\externaldocument[stacks-limits-]{stacks-limits}
\externaldocument[stacks-cohomology-]{stacks-cohomology}
\externaldocument[stacks-perfect-]{stacks-perfect}
\externaldocument[stacks-introduction-]{stacks-introduction}
\externaldocument[stacks-more-morphisms-]{stacks-more-morphisms}
\externaldocument[stacks-geometry-]{stacks-geometry}
\externaldocument[moduli-]{moduli}
\externaldocument[moduli-curves-]{moduli-curves}
\externaldocument[examples-]{examples}
\externaldocument[exercises-]{exercises}
\externaldocument[guide-]{guide}
\externaldocument[desirables-]{desirables}
\externaldocument[coding-]{coding}
\externaldocument[obsolete-]{obsolete}
\externaldocument[fdl-]{fdl}
\externaldocument[index-]{index}

% Theorem environments.
%
\theoremstyle{plain}
\newtheorem{theorem}[subsection]{Theorem}
\newtheorem{proposition}[subsection]{Proposition}
\newtheorem{lemma}[subsection]{Lemma}

\theoremstyle{definition}
\newtheorem{definition}[subsection]{Definition}
\newtheorem{example}[subsection]{Example}
\newtheorem{exercise}[subsection]{Exercise}
\newtheorem{situation}[subsection]{Situation}

\theoremstyle{remark}
\newtheorem{remark}[subsection]{Remark}
\newtheorem{remarks}[subsection]{Remarks}

\numberwithin{equation}{subsection}

% Macros
%
\def\lim{\mathop{\mathrm{lim}}\nolimits}
\def\colim{\mathop{\mathrm{colim}}\nolimits}
\def\Spec{\mathop{\mathrm{Spec}}}
\def\Hom{\mathop{\mathrm{Hom}}\nolimits}
\def\Ext{\mathop{\mathrm{Ext}}\nolimits}
\def\SheafHom{\mathop{\mathcal{H}\!\mathit{om}}\nolimits}
\def\SheafExt{\mathop{\mathcal{E}\!\mathit{xt}}\nolimits}
\def\Sch{\mathit{Sch}}
\def\Mor{\mathop{\mathrm{Mor}}\nolimits}
\def\Ob{\mathop{\mathrm{Ob}}\nolimits}
\def\Sh{\mathop{\mathit{Sh}}\nolimits}
\def\NL{\mathop{N\!L}\nolimits}
\def\CH{\mathop{\mathrm{CH}}\nolimits}
\def\proetale{{pro\text{-}\acute{e}tale}}
\def\etale{{\acute{e}tale}}
\def\QCoh{\mathit{QCoh}}
\def\Ker{\mathop{\mathrm{Ker}}}
\def\Im{\mathop{\mathrm{Im}}}
\def\Coker{\mathop{\mathrm{Coker}}}
\def\Coim{\mathop{\mathrm{Coim}}}

% Boxtimes
%
\DeclareMathSymbol{\boxtimes}{\mathbin}{AMSa}{"02}

%
% Macros for moduli stacks/spaces
%
\def\QCohstack{\mathcal{QC}\!\mathit{oh}}
\def\Cohstack{\mathcal{C}\!\mathit{oh}}
\def\Spacesstack{\mathcal{S}\!\mathit{paces}}
\def\Quotfunctor{\mathrm{Quot}}
\def\Hilbfunctor{\mathrm{Hilb}}
\def\Curvesstack{\mathcal{C}\!\mathit{urves}}
\def\Polarizedstack{\mathcal{P}\!\mathit{olarized}}
\def\Complexesstack{\mathcal{C}\!\mathit{omplexes}}
% \Pic is the operator that assigns to X its picard group, usage \Pic(X)
% \Picardstack_{X/B} denotes the Picard stack of X over B
% \Picardfunctor_{X/B} denotes the Picard functor of X over B
\def\Pic{\mathop{\mathrm{Pic}}\nolimits}
\def\Picardstack{\mathcal{P}\!\mathit{ic}}
\def\Picardfunctor{\mathrm{Pic}}
\def\Deformationcategory{\mathcal{D}\!\mathit{ef}}
\usepackage{cleveref}
\crefname{originallemma}{original lemma}{original lemmas}
\Crefname{originallemma}{Original lemma}{Original lemmas}
\makeatletter
\newlabel{derived-lemma-right-adjoint}{{0CQS}{source}{}{sourcepointer-1}{}}
\newlabel{derived-lemma-right-adjoint@cref}{{[originallemma][0][]0CQS}{[1][0][]source}{}{}{}}
\makeatother
\crefname{originallemma}{original lemma}{original lemmas}
\Crefname{originallemma}{Original lemma}{Original lemmas}
\makeatletter
\newlabel{derived-lemma-prepare-adjoint}{{0CQR}{source}{}{sourcepointer-2}{}}
\newlabel{derived-lemma-prepare-adjoint@cref}{{[originallemma][0][]0CQR}{[1][0][]source}{}{}{}}
\makeatother
\crefname{originallemma}{original lemma}{original lemmas}
\Crefname{originallemma}{Original lemma}{Original lemmas}
\makeatletter
\newlabel{coherent-lemma-induction-principle}{{08DR}{source}{}{sourcepointer-3}{}}
\newlabel{coherent-lemma-induction-principle@cref}{{[originallemma][0][]08DR}{[1][0][]source}{}{}{}}
\makeatother
\crefname{originalremark}{original remark}{original remarks}
\Crefname{originalremark}{Original remark}{Original remarks}
\makeatletter
\newlabel{section-coherator}{{08D6}{source}{}{sourcepointer-4}{}}
\newlabel{section-coherator@cref}{{[originalremark][0][]08D6}{[1][0][]source}{}{}{}}
\makeatother
\crefname{originallemma}{original lemma}{original lemmas}
\Crefname{originallemma}{Original lemma}{Original lemmas}
\makeatletter
\newlabel{lemma-affine-coherator}{{08D9}{source}{}{sourcepointer-5}{}}
\newlabel{lemma-affine-coherator@cref}{{[originallemma][0][]08D9}{[1][0][]source}{}{}{}}
\makeatother
\crefname{originallemma}{original lemma}{original lemmas}
\Crefname{originallemma}{Original lemma}{Original lemmas}
\makeatletter
\newlabel{cohomology-lemma-exact-sequence-j-star}{{08BV}{source}{}{sourcepointer-6}{}}
\newlabel{cohomology-lemma-exact-sequence-j-star@cref}{{[originallemma][0][]08BV}{[1][0][]source}{}{}{}}
\makeatother
\crefname{originallemma}{original lemma}{original lemmas}
\Crefname{originallemma}{Original lemma}{Original lemmas}
\makeatletter
\newlabel{lemma-quasi-coherence-direct-image}{{08D5}{source}{}{sourcepointer-7}{}}
\newlabel{lemma-quasi-coherence-direct-image@cref}{{[originallemma][0][]08D5}{[1][0][]source}{}{}{}}
\makeatother
\crefname{originallemma}{original lemma}{original lemmas}
\Crefname{originallemma}{Original lemma}{Original lemmas}
\makeatletter
\newlabel{lemma-quasi-coherence-pullback}{{08DW}{source}{}{sourcepointer-8}{}}
\newlabel{lemma-quasi-coherence-pullback@cref}{{[originallemma][0][]08DW}{[1][0][]source}{}{}{}}
\makeatother
\title{The inclusion of the quasi-coherent derived category has a right adjoint}
\author{The Stacks Project Authors}
\date{stacks-project:perfect.tex:local-snapshot-20261012; GNU FDL 1.2 or later}
\begin{document}
\maketitle
\setcounter{subsection}{1}
\noindent\textit{Editorial note.} All delivered mathematical text is the author's own, copied from the pinned Stacks Project chapter file perfect.tex; no statement or proof step was rewritten, reordered or removed. The counted span is the author environment and its own proof environment. References inside the span are handed to the assembler as navigation stubs that carry the canonical tag of the cited tag only; whether each of them is load-bearing is recorded on the stub itself.\par\medskip
\noindent\textit{Allowed standard prerequisites (editorial).}\begin{itemize}
\item A functor has a right adjoint precisely when every object sits in a distinguished triangle with a quasi-coherent first term and a last term orthogonal to all quasi-coherent objects. invoked at proof lines 5051-5080. Reduces the construction of the adjoint to producing such a triangle for every object.
\item On an affine scheme the inclusion of the quasi-coherent derived category has the right adjoint given by the coherator. invoked at proof lines 5055-5062. Provides the base case of the induction.
\item The induction principle on quasi-compact opens reduces the statement about a quasi-compact quasi-separated scheme to the base case and to a union of two quasi-compact opens. invoked at proof lines 5055-5063. Organises the proof into the base case and the union step.
\item Pushforwards and pullbacks along open immersions preserve the quasi-coherent derived category, and the standard exact sequence relates the pushforwards from two opens and their intersection. invoked at proof lines 5081-5119. Produces the required triangle in the union step and the orthogonality computation.
\end{itemize}\medskip
\ifdefined\hypertarget\hypertarget{sourcepointer-1}{}\fi
\noindent\textit{Source reference.} Original lemma 0CQS: sources/stacks-chow/tags.txt, source lines 1--1. Invoked at proof lines 5051-5060 for the reduction to the triangle criterion and again at proof lines 5090-5095 in the if and only if form; both uses are inside the counted span.\par\medskip
\ifdefined\hypertarget\hypertarget{sourcepointer-2}{}\fi
\noindent\textit{Source reference.} Original lemma 0CQR: sources/stacks-chow/tags.txt, source lines 1--1. A navigation pointer in the adjoint construction; the construction carried out with it is inside the counted span.\par\medskip
\ifdefined\hypertarget\hypertarget{sourcepointer-3}{}\fi
\noindent\textit{Source reference.} Original lemma 08DR: sources/stacks-chow/tags.txt, source lines 1--1. Cited at proof lines 5055-5060 to organise the proof into the base case and the union step; both are inside the counted span.\par\medskip
\ifdefined\hypertarget\hypertarget{sourcepointer-4}{}\fi
\noindent\textit{Source reference.} Original remark 08D6: sources/stacks-perfect/pinned-original-source.tex, source lines 1543--1543. Cited at proof lines 5057-5060 for the existence of the right derived functor used in the base case; the base case argument is inside the counted span.\par\medskip
\ifdefined\hypertarget\hypertarget{sourcepointer-5}{}\fi
\noindent\textit{Source reference.} Original lemma 08D9: sources/stacks-perfect/pinned-original-source.tex, source lines 1670--1670. Cited at proof lines 5061-5063 to settle the affine case; the induction set up that uses the base case is inside the counted span.\par\medskip
\ifdefined\hypertarget\hypertarget{sourcepointer-6}{}\fi
\noindent\textit{Source reference.} Original lemma 08BV: sources/stacks-chow/tags.txt, source lines 1--1. Used at proof lines 5071-5080 for the displayed triangle in the union step; the verification of its properties is inside the counted span.\par\medskip
\ifdefined\hypertarget\hypertarget{sourcepointer-7}{}\fi
\noindent\textit{Source reference.} Original lemma 08D5: sources/stacks-perfect/pinned-original-source.tex, source lines 661--661. Invoked at proof lines 5106-5110 to see that the pushed forward quasi-coherent term is quasi-coherent; the second required property is verified immediately after inside the counted span.\par\medskip
\ifdefined\hypertarget\hypertarget{sourcepointer-8}{}\fi
\noindent\textit{Source reference.} Original lemma 08DW: sources/stacks-perfect/pinned-original-source.tex, source lines 581--581. Used at proof lines 5112-5118 to see that the pullback of a quasi-coherent object stays quasi-coherent, which combined with adjointness gives the vanishing; that computation is inside the counted span.\par\medskip
\begin{lemma}
\label{lemma-better-coherator}
Let $X$ be a quasi-compact and quasi-separated scheme.
The inclusion functor $D_\QCoh(\mathcal{O}_X) \to D(\mathcal{O}_X)$
has a right adjoint $DQ_X$.
\end{lemma}

\begin{proof}[First proof]
We will use the induction principle as in
Cohomology of Schemes, Lemma \ref{coherent-lemma-induction-principle}
to prove this. If $D(\QCoh(\mathcal{O}_X)) \to D_\QCoh(\mathcal{O}_X)$
is an equivalence, then the lemma is true because the functor
$RQ_X$ of Section \ref{section-coherator} is a right adjoint to the functor
$D(\QCoh(\mathcal{O}_X)) \to D(\mathcal{O}_X)$.
In particular, our lemma is true for affine schemes, see
Lemma \ref{lemma-affine-coherator}.
Thus we see that it suffices to show: if $X = U \cup V$
is a union of two quasi-compact opens and the
lemma holds for $U$, $V$, and $U \cap V$, then the lemma holds for $X$.

\medskip\noindent
The adjoint exists if and only if for every object $K$ of
$D(\mathcal{O}_X)$ we can find a distinguished triangle
$$
E' \to E \to K \to E'[1]
$$
in $D(\mathcal{O}_X)$
such that $E'$ is in $D_\QCoh(\mathcal{O}_X)$ and such that
$\Hom(M, K) = 0$ for all $M$ in $D_\QCoh(\mathcal{O}_X)$. See
Derived Categories, Lemma \ref{derived-lemma-right-adjoint}.
Consider the distinguished triangle
$$
E \to Rj_{U, *}E|_U \oplus Rj_{V, *}E|_V \to
Rj_{U \cap V, *}E|_{U \cap V} \to E[1]
$$
in $D(\mathcal{O}_X)$ of
Cohomology, Lemma \ref{cohomology-lemma-exact-sequence-j-star}.
By Derived Categories, Lemma \ref{derived-lemma-prepare-adjoint}
it suffices to construct the desired distinguished triangles
for $Rj_{U, *}E|_U$, $Rj_{V, *}E|_V$, and
$Rj_{U \cap V, *}E|_{U \cap V}$. This reduces us to the statement
discussed in the next paragraph.

\medskip\noindent
Let $j : U \to X$ be an open immersion corresponding with $U$ a quasi-compact
open for which the lemma is true. Let $L$ be an object of $D(\mathcal{O}_U)$.
Then there exists a distinguished triangle
$$
E' \to Rj_*L \to K \to E'[1]
$$
in $D(\mathcal{O}_X)$
such that $E'$ is in $D_\QCoh(\mathcal{O}_X)$ and such that
$\Hom(M, K) = 0$ for all $M$ in $D_\QCoh(\mathcal{O}_X)$.
To see this we choose a distinguished triangle
$$
L' \to L \to Q \to L'[1]
$$
in $D(\mathcal{O}_U)$ such that $L'$ is in $D_\QCoh(\mathcal{O}_U)$
and such that $\Hom(N, Q) = 0$ for all $N$ in $D_\QCoh(\mathcal{O}_U)$.
This is possible because the statement in
Derived Categories, Lemma \ref{derived-lemma-right-adjoint}
is an if and only if.
We obtain a distinguished triangle
$$
Rj_*L' \to Rj_*L \to Rj_*Q \to Rj_*L'[1]
$$
in $D(\mathcal{O}_X)$. Observe that $Rj_*L'$ is in $D_\QCoh(\mathcal{O}_X)$
by Lemma \ref{lemma-quasi-coherence-direct-image}.
On the other hand, if $M$ in $D_\QCoh(\mathcal{O}_X)$, then
$$
\Hom(M, Rj_*Q) = \Hom(Lj^*M, Q) = 0
$$
because $Lj^*M$ is in $D_\QCoh(\mathcal{O}_U)$ by
Lemma \ref{lemma-quasi-coherence-pullback}.
This finishes the proof.
\end{proof}

\end{document}
